% BHU PYQ -- Transcribed & Corrected
% Paper: BOTMJ-45 / BOTMJ45: Plant Ecology
% Exam: B.Sc. (Hons) Semester IV Examination, 2025-26 (NEP)
% Category: Major Core
% Department: Department of Botany, Institute of Science / Faculty of Science, Banaras Hindu University

\documentclass[11pt,a4paper]{article}
\usepackage[a4paper,top=2.5cm,bottom=2.5cm,left=2.8cm,right=2.8cm,headheight=14pt]{geometry}
\usepackage{amsmath,amssymb}
\usepackage[T1]{fontenc}
\usepackage[utf8]{inputenc}
\usepackage{lmodern,microtype,fancyhdr,enumitem,hyperref,lastpage,parskip}

\hypersetup{
    colorlinks=true,
    urlcolor=black,
    linkcolor=black,
    pdftitle={BOTMJ45: Plant Ecology -- BHU B.Sc. Sem IV 2025-26 (NEP)}
}

\pagestyle{fancy}
\fancyhf{}
\renewcommand{\headrulewidth}{0.4pt}
\renewcommand{\footrulewidth}{0.4pt}
\fancyhead[L]{\small \textit{Banaras Hindu University}}
\fancyhead[R]{\small \textit{BOTMJ-45: Plant Ecology (Sem IV 2025-26)}}
\fancyfoot[C]{\small Page \thepage\ of \pageref{LastPage}}

\newlist{parts}{enumerate}{2}
\setlist[parts,1]{label=(\alph*),leftmargin=*,itemsep=4pt,topsep=2pt}

\newcommand{\pts}[1]{\hfill \textit{[#1]}}

\begin{document}

\begin{center}
  {\large\bfseries BANARAS HINDU UNIVERSITY}\\[2pt]
  {\normalsize B.Sc. (Hons) Semester IV Examination, 2025-26 (NEP)}\\[6pt]
  \rule{0.8\linewidth}{0.5pt}\\[6pt]
  {\large\bfseries BOTANY}\\[4pt]
  {\large\bfseries Paper No. BOTMJ45 : Plant Ecology}\\[2pt]
  {\small Ref. No.: Conf/Even/N/2025-26/21}\\[4pt]
  {\normalsize Time: $2\frac{1}{2}$ hours \hfill Full Marks: 70}
\end{center}

\rule{\linewidth}{0.4pt}

\smallskip

\noindent \textit{(Write your Roll No. at the top immediately on the receipt of this question paper)}\\[2pt]
\noindent \textit{The figures in the right-hand margin indicate marks.}\\[1pt]
\noindent \textit{Answer \textbf{five} questions including Question 1 which is compulsory.}

\bigskip

\noindent \textbf{Question 1.} Describe briefly/define/answer, in about 50 words each of the following: \pts{7 $\times$ 2 = 14}
\begin{parts}
  \item Synthetic community characters \pts{2}
  \item Soil forming factors \pts{2}
  \item Define phenology. \pts{2}
  \item Raunkiaer's life forms \pts{2}
  \item Distinguish between phototrophs and chemotrophs \pts{2}
  \item Function of stomatal hairs in plant adaptation \pts{2}
  \item What is ecesis? \pts{2}
\end{parts}

\bigskip
\rule{\linewidth}{0.2pt}
\bigskip

\noindent \textbf{Question 2.} Describe soil structure and explain soil profile with the help of a suitable diagram. \pts{14}

\medskip\rule{\linewidth}{0.2pt}\medskip

\noindent \textbf{Question 3.} Write short notes on \textbf{any two} of the following: \pts{2 $\times$ 7 = 14}
\begin{parts}
  \item Atmosphere \pts{7}
  \item Adaptations in hydrophytes \pts{7}
  \item Role of temperature in distribution of vegetation \pts{7}
\end{parts}

\medskip\rule{\linewidth}{0.2pt}\medskip

\noindent \textbf{Question 4.} Why biotic interaction is important in community? Describe various types of positive and negative biotic interactions in a community. \pts{14}

\medskip\rule{\linewidth}{0.2pt}\medskip

\noindent \textbf{Question 5.} Write short notes on \textbf{any two} of the following: \pts{2 $\times$ 7 = 14}
\begin{parts}
  \item Survivorship curve \pts{7}
  \item Biological spectrum \pts{7}
  \item Causes of ecological succession \pts{7}
\end{parts}

\medskip\rule{\linewidth}{0.2pt}\medskip

\noindent \textbf{Question 6.} Define community. Discuss various quantitative characters of the community. \pts{14}

\medskip\rule{\linewidth}{0.2pt}\medskip

\noindent \textbf{Question 7.} Write short notes on \textbf{any two} of the following: \pts{2 $\times$ 7 = 14}
\begin{parts}
  \item Nitrogen cycle \pts{7}
  \item Models of energy flow in the ecosystem \pts{7}
  \item Ecological pyramids \pts{7}
\end{parts}

\vfill
\rule{\linewidth}{0.4pt}

\begin{center}\small
\href{https://bhu-exam.vercel.app/}{Click here to visit our website}\\[4pt]
\href{https://me-aryan.vercel.app/}{Click here to visit our contributor}
\end{center}

\end{document}
